% Tabelle der offenen Positionen am Desktop (Reiter Active, blaues Design) – Vektor-Nachbau
% als AUSSCHNITT mit Auslassung in der Mitte (Lehrbuch-Darstellung), damit die Schrift bei
% 170 mm Breite nicht unter 6 pt fällt.
% Grundlage: DOM-Messung originals/messungen/positionen-desktop.json (CSS-px, Fenster
% 2040 x 986; anonymisiert: Ziffern = 0, Instrumente = XXXXXX). Koordinaten unten in
% CSS-px, y nach unten, Ursprung = linke obere Ecke des Ausschnitts.
%
% Gemessen (Original-x):
%   .maintable   1704 x 271, Fläche rgb(4,39,60), Radius 7 px, Inter 400, 12 px / 14 px
%   ul.nav-tabs  x1436,3 y10 257,7 x 30, Fläche rgb(0,19,40), Radius 3 px; endet bündig
%                mit der Tabelle (x1694 = 1704 - 10)
%     a.nav-link  flex, padding 0 16px, Text mittig, weiß, capitalize:
%                „Active (0)“ 87 (aktiv: Fläche rgb(14,80,114), Radius 3 0 0 3),
%                „Pending (0)“ 98,1, „History“ 72,6 (Radius 0 3 3 0)
%   th           y47, 28,8 hoch, padding 4px 10px, Text rgb(168,187,190), capitalize;
%                span (inline-block, 20 px Zeile) mit padding-right 15 px; Text beginnt am
%                gemessenen span-x: Instrument 177, Type 298, Amount 383,8,
%                Trade volume 493,5, Swap 1113,6, Commission 1207,4, Total Profit 1344,2,
%                Margin 1467,9
%   td           Zeilen y75,8 / y105,4, je 29,6 hoch, border-top 1,6 px rgb(4,39,60),
%                Fläche rgba(14,80,114,.1), Text mittig (padding 4px 10px); erste Zelle
%                Radius 3 0 0 3, letzte 0 3 3 0
%                Spalten: Instrument x151,8 w126,2 | Type 278 w83,2 | Amount 361,2 w104,3 |
%                Trade volume 465,5 w149,5 | Swap 1093 w87,8 | Commission 1180,8 w138,1 |
%                Total Profit 1318,9 w127,2 | Margin 1446 w98
%   Total Profit grün rgb(83,188,81) (color-green-main) bzw. rot rgb(249,64,86)
%                (color-red-main)
% Grundlinien (Chrome, Gerätepixel 0,8 px: Ascent 12, Descent 3,2 px): Zeilenoberkante +
% (Zeilenhöhe - 15,2) / 2 + 12. Belegt durch die Messung: €-span der Zeile 1 oben bei 83
% (+ 12 = 95,0), Mitte von Bild und Pfeil 91,7 (+ halbe x-Höhe 3,3 = 95,0).
%   Kopf: 51 + 2,4 + 12 = 65,4; Zeilen: 95,0 und 124,6; Reiter (14 px Zeile mittig in
%   30 px): 18 - 0,8 + 12 = 29,2 (mittige Ausrichtung im Reiter angenommen, nicht gemessen)
%
% Ausschnitt (Abweichungen vom Original, nur Weglassen und Zusammenschieben):
%   links  Spalten Instrument bis Trade volume (Original-x 151,8..615), um -141,8 verschoben
%          (Bild- und ID-Spalte entfallen; die Zeilen beginnen wieder mit Radius 3 px);
%   Mitte  Auslassung als Zickzack-Bruchkante über die volle Höhe (Lücke 10 px, Fläche frei):
%          Open Rate, Open Time, S/L, T/P;
%   rechts Spalten Swap bis Margin (Original-x 1093..1544), um -597,8 verschoben
%          (Zahnrad- und Schließen-Spalte, 150 px, entfallen; Zeilenende wieder Radius 3 px).
%          Die Reiter bleiben bündig mit dem rechten Tabellenrand.
%   unten  Schnitt 10 px unter der zweiten Zeile (y145), Ecken wieder Radius 7 px.
%   Weggelassen: Leiste New order / Search / All sides / All profits (liegt links außerhalb
%   des Ausschnitts), Griff zum Einklappen (a.hide, Mitte oben), Instrument-Bild und Pfeil
%   hinter dem Instrument (PNG, kein Vektor; die Fläche bleibt leer: Text um 5,5 px nach
%   links wie im Original), Sortierpfeile (padding-right 15 px der Köpfe; Inhalt nicht
%   gemessen, nicht gezeichnet).
%
% Das Dokument ist zahlenfrei: Werte sind Wörter, ohne „€“ (wie bar.tex), Commission „—“.
% Reiter ohne Zähler; ihre Breite rechnet wie der Browser (16 + Text + 16 px). Probe der
% Textbreiten mit den gemessenen Werten: siehe \typeout POSITIONEN-DESKTOP im Log.
%
% Marken (Schalter \ifmarken, Standard: an; ohne Marken bauen mit
%   xelatex '\def\ohnemarken{}% Tabelle der offenen Positionen am Desktop (Reiter Active, blaues Design) – Vektor-Nachbau
% als AUSSCHNITT mit Auslassung in der Mitte (Lehrbuch-Darstellung), damit die Schrift bei
% 170 mm Breite nicht unter 6 pt fällt.
% Grundlage: DOM-Messung originals/messungen/positionen-desktop.json (CSS-px, Fenster
% 2040 x 986; anonymisiert: Ziffern = 0, Instrumente = XXXXXX). Koordinaten unten in
% CSS-px, y nach unten, Ursprung = linke obere Ecke des Ausschnitts.
%
% Gemessen (Original-x):
%   .maintable   1704 x 271, Fläche rgb(4,39,60), Radius 7 px, Inter 400, 12 px / 14 px
%   ul.nav-tabs  x1436,3 y10 257,7 x 30, Fläche rgb(0,19,40), Radius 3 px; endet bündig
%                mit der Tabelle (x1694 = 1704 - 10)
%     a.nav-link  flex, padding 0 16px, Text mittig, weiß, capitalize:
%                „Active (0)“ 87 (aktiv: Fläche rgb(14,80,114), Radius 3 0 0 3),
%                „Pending (0)“ 98,1, „History“ 72,6 (Radius 0 3 3 0)
%   th           y47, 28,8 hoch, padding 4px 10px, Text rgb(168,187,190), capitalize;
%                span (inline-block, 20 px Zeile) mit padding-right 15 px; Text beginnt am
%                gemessenen span-x: Instrument 177, Type 298, Amount 383,8,
%                Trade volume 493,5, Swap 1113,6, Commission 1207,4, Total Profit 1344,2,
%                Margin 1467,9
%   td           Zeilen y75,8 / y105,4, je 29,6 hoch, border-top 1,6 px rgb(4,39,60),
%                Fläche rgba(14,80,114,.1), Text mittig (padding 4px 10px); erste Zelle
%                Radius 3 0 0 3, letzte 0 3 3 0
%                Spalten: Instrument x151,8 w126,2 | Type 278 w83,2 | Amount 361,2 w104,3 |
%                Trade volume 465,5 w149,5 | Swap 1093 w87,8 | Commission 1180,8 w138,1 |
%                Total Profit 1318,9 w127,2 | Margin 1446 w98
%   Total Profit grün rgb(83,188,81) (color-green-main) bzw. rot rgb(249,64,86)
%                (color-red-main)
% Grundlinien (Chrome, Gerätepixel 0,8 px: Ascent 12, Descent 3,2 px): Zeilenoberkante +
% (Zeilenhöhe - 15,2) / 2 + 12. Belegt durch die Messung: €-span der Zeile 1 oben bei 83
% (+ 12 = 95,0), Mitte von Bild und Pfeil 91,7 (+ halbe x-Höhe 3,3 = 95,0).
%   Kopf: 51 + 2,4 + 12 = 65,4; Zeilen: 95,0 und 124,6; Reiter (14 px Zeile mittig in
%   30 px): 18 - 0,8 + 12 = 29,2 (mittige Ausrichtung im Reiter angenommen, nicht gemessen)
%
% Ausschnitt (Abweichungen vom Original, nur Weglassen und Zusammenschieben):
%   links  Spalten Instrument bis Trade volume (Original-x 151,8..615), um -141,8 verschoben
%          (Bild- und ID-Spalte entfallen; die Zeilen beginnen wieder mit Radius 3 px);
%   Mitte  Auslassung als Zickzack-Bruchkante über die volle Höhe (Lücke 10 px, Fläche frei):
%          Open Rate, Open Time, S/L, T/P;
%   rechts Spalten Swap bis Margin (Original-x 1093..1544), um -597,8 verschoben
%          (Zahnrad- und Schließen-Spalte, 150 px, entfallen; Zeilenende wieder Radius 3 px).
%          Die Reiter bleiben bündig mit dem rechten Tabellenrand.
%   unten  Schnitt 10 px unter der zweiten Zeile (y145), Ecken wieder Radius 7 px.
%   Weggelassen: Leiste New order / Search / All sides / All profits (liegt links außerhalb
%   des Ausschnitts), Griff zum Einklappen (a.hide, Mitte oben), Instrument-Bild und Pfeil
%   hinter dem Instrument (PNG, kein Vektor; die Fläche bleibt leer: Text um 5,5 px nach
%   links wie im Original), Sortierpfeile (padding-right 15 px der Köpfe; Inhalt nicht
%   gemessen, nicht gezeichnet).
%
% Das Dokument ist zahlenfrei: Werte sind Wörter, ohne „€“ (wie bar.tex), Commission „—“.
% Reiter ohne Zähler; ihre Breite rechnet wie der Browser (16 + Text + 16 px). Probe der
% Textbreiten mit den gemessenen Werten: siehe \typeout POSITIONEN-DESKTOP im Log.
%
% Marken (Schalter \ifmarken, Standard: an; ohne Marken bauen mit
%   xelatex '\def\ohnemarken{}% Tabelle der offenen Positionen am Desktop (Reiter Active, blaues Design) – Vektor-Nachbau
% als AUSSCHNITT mit Auslassung in der Mitte (Lehrbuch-Darstellung), damit die Schrift bei
% 170 mm Breite nicht unter 6 pt fällt.
% Grundlage: DOM-Messung originals/messungen/positionen-desktop.json (CSS-px, Fenster
% 2040 x 986; anonymisiert: Ziffern = 0, Instrumente = XXXXXX). Koordinaten unten in
% CSS-px, y nach unten, Ursprung = linke obere Ecke des Ausschnitts.
%
% Gemessen (Original-x):
%   .maintable   1704 x 271, Fläche rgb(4,39,60), Radius 7 px, Inter 400, 12 px / 14 px
%   ul.nav-tabs  x1436,3 y10 257,7 x 30, Fläche rgb(0,19,40), Radius 3 px; endet bündig
%                mit der Tabelle (x1694 = 1704 - 10)
%     a.nav-link  flex, padding 0 16px, Text mittig, weiß, capitalize:
%                „Active (0)“ 87 (aktiv: Fläche rgb(14,80,114), Radius 3 0 0 3),
%                „Pending (0)“ 98,1, „History“ 72,6 (Radius 0 3 3 0)
%   th           y47, 28,8 hoch, padding 4px 10px, Text rgb(168,187,190), capitalize;
%                span (inline-block, 20 px Zeile) mit padding-right 15 px; Text beginnt am
%                gemessenen span-x: Instrument 177, Type 298, Amount 383,8,
%                Trade volume 493,5, Swap 1113,6, Commission 1207,4, Total Profit 1344,2,
%                Margin 1467,9
%   td           Zeilen y75,8 / y105,4, je 29,6 hoch, border-top 1,6 px rgb(4,39,60),
%                Fläche rgba(14,80,114,.1), Text mittig (padding 4px 10px); erste Zelle
%                Radius 3 0 0 3, letzte 0 3 3 0
%                Spalten: Instrument x151,8 w126,2 | Type 278 w83,2 | Amount 361,2 w104,3 |
%                Trade volume 465,5 w149,5 | Swap 1093 w87,8 | Commission 1180,8 w138,1 |
%                Total Profit 1318,9 w127,2 | Margin 1446 w98
%   Total Profit grün rgb(83,188,81) (color-green-main) bzw. rot rgb(249,64,86)
%                (color-red-main)
% Grundlinien (Chrome, Gerätepixel 0,8 px: Ascent 12, Descent 3,2 px): Zeilenoberkante +
% (Zeilenhöhe - 15,2) / 2 + 12. Belegt durch die Messung: €-span der Zeile 1 oben bei 83
% (+ 12 = 95,0), Mitte von Bild und Pfeil 91,7 (+ halbe x-Höhe 3,3 = 95,0).
%   Kopf: 51 + 2,4 + 12 = 65,4; Zeilen: 95,0 und 124,6; Reiter (14 px Zeile mittig in
%   30 px): 18 - 0,8 + 12 = 29,2 (mittige Ausrichtung im Reiter angenommen, nicht gemessen)
%
% Ausschnitt (Abweichungen vom Original, nur Weglassen und Zusammenschieben):
%   links  Spalten Instrument bis Trade volume (Original-x 151,8..615), um -141,8 verschoben
%          (Bild- und ID-Spalte entfallen; die Zeilen beginnen wieder mit Radius 3 px);
%   Mitte  Auslassung als Zickzack-Bruchkante über die volle Höhe (Lücke 10 px, Fläche frei):
%          Open Rate, Open Time, S/L, T/P;
%   rechts Spalten Swap bis Margin (Original-x 1093..1544), um -597,8 verschoben
%          (Zahnrad- und Schließen-Spalte, 150 px, entfallen; Zeilenende wieder Radius 3 px).
%          Die Reiter bleiben bündig mit dem rechten Tabellenrand.
%   unten  Schnitt 10 px unter der zweiten Zeile (y145), Ecken wieder Radius 7 px.
%   Weggelassen: Leiste New order / Search / All sides / All profits (liegt links außerhalb
%   des Ausschnitts), Griff zum Einklappen (a.hide, Mitte oben), Instrument-Bild und Pfeil
%   hinter dem Instrument (PNG, kein Vektor; die Fläche bleibt leer: Text um 5,5 px nach
%   links wie im Original), Sortierpfeile (padding-right 15 px der Köpfe; Inhalt nicht
%   gemessen, nicht gezeichnet).
%
% Das Dokument ist zahlenfrei: Werte sind Wörter, ohne „€“ (wie bar.tex), Commission „—“.
% Reiter ohne Zähler; ihre Breite rechnet wie der Browser (16 + Text + 16 px). Probe der
% Textbreiten mit den gemessenen Werten: siehe \typeout POSITIONEN-DESKTOP im Log.
%
% Marken (Schalter \ifmarken, Standard: an; ohne Marken bauen mit
%   xelatex '\def\ohnemarken{}% Tabelle der offenen Positionen am Desktop (Reiter Active, blaues Design) – Vektor-Nachbau
% als AUSSCHNITT mit Auslassung in der Mitte (Lehrbuch-Darstellung), damit die Schrift bei
% 170 mm Breite nicht unter 6 pt fällt.
% Grundlage: DOM-Messung originals/messungen/positionen-desktop.json (CSS-px, Fenster
% 2040 x 986; anonymisiert: Ziffern = 0, Instrumente = XXXXXX). Koordinaten unten in
% CSS-px, y nach unten, Ursprung = linke obere Ecke des Ausschnitts.
%
% Gemessen (Original-x):
%   .maintable   1704 x 271, Fläche rgb(4,39,60), Radius 7 px, Inter 400, 12 px / 14 px
%   ul.nav-tabs  x1436,3 y10 257,7 x 30, Fläche rgb(0,19,40), Radius 3 px; endet bündig
%                mit der Tabelle (x1694 = 1704 - 10)
%     a.nav-link  flex, padding 0 16px, Text mittig, weiß, capitalize:
%                „Active (0)“ 87 (aktiv: Fläche rgb(14,80,114), Radius 3 0 0 3),
%                „Pending (0)“ 98,1, „History“ 72,6 (Radius 0 3 3 0)
%   th           y47, 28,8 hoch, padding 4px 10px, Text rgb(168,187,190), capitalize;
%                span (inline-block, 20 px Zeile) mit padding-right 15 px; Text beginnt am
%                gemessenen span-x: Instrument 177, Type 298, Amount 383,8,
%                Trade volume 493,5, Swap 1113,6, Commission 1207,4, Total Profit 1344,2,
%                Margin 1467,9
%   td           Zeilen y75,8 / y105,4, je 29,6 hoch, border-top 1,6 px rgb(4,39,60),
%                Fläche rgba(14,80,114,.1), Text mittig (padding 4px 10px); erste Zelle
%                Radius 3 0 0 3, letzte 0 3 3 0
%                Spalten: Instrument x151,8 w126,2 | Type 278 w83,2 | Amount 361,2 w104,3 |
%                Trade volume 465,5 w149,5 | Swap 1093 w87,8 | Commission 1180,8 w138,1 |
%                Total Profit 1318,9 w127,2 | Margin 1446 w98
%   Total Profit grün rgb(83,188,81) (color-green-main) bzw. rot rgb(249,64,86)
%                (color-red-main)
% Grundlinien (Chrome, Gerätepixel 0,8 px: Ascent 12, Descent 3,2 px): Zeilenoberkante +
% (Zeilenhöhe - 15,2) / 2 + 12. Belegt durch die Messung: €-span der Zeile 1 oben bei 83
% (+ 12 = 95,0), Mitte von Bild und Pfeil 91,7 (+ halbe x-Höhe 3,3 = 95,0).
%   Kopf: 51 + 2,4 + 12 = 65,4; Zeilen: 95,0 und 124,6; Reiter (14 px Zeile mittig in
%   30 px): 18 - 0,8 + 12 = 29,2 (mittige Ausrichtung im Reiter angenommen, nicht gemessen)
%
% Ausschnitt (Abweichungen vom Original, nur Weglassen und Zusammenschieben):
%   links  Spalten Instrument bis Trade volume (Original-x 151,8..615), um -141,8 verschoben
%          (Bild- und ID-Spalte entfallen; die Zeilen beginnen wieder mit Radius 3 px);
%   Mitte  Auslassung als Zickzack-Bruchkante über die volle Höhe (Lücke 10 px, Fläche frei):
%          Open Rate, Open Time, S/L, T/P;
%   rechts Spalten Swap bis Margin (Original-x 1093..1544), um -597,8 verschoben
%          (Zahnrad- und Schließen-Spalte, 150 px, entfallen; Zeilenende wieder Radius 3 px).
%          Die Reiter bleiben bündig mit dem rechten Tabellenrand.
%   unten  Schnitt 10 px unter der zweiten Zeile (y145), Ecken wieder Radius 7 px.
%   Weggelassen: Leiste New order / Search / All sides / All profits (liegt links außerhalb
%   des Ausschnitts), Griff zum Einklappen (a.hide, Mitte oben), Instrument-Bild und Pfeil
%   hinter dem Instrument (PNG, kein Vektor; die Fläche bleibt leer: Text um 5,5 px nach
%   links wie im Original), Sortierpfeile (padding-right 15 px der Köpfe; Inhalt nicht
%   gemessen, nicht gezeichnet).
%
% Das Dokument ist zahlenfrei: Werte sind Wörter, ohne „€“ (wie bar.tex), Commission „—“.
% Reiter ohne Zähler; ihre Breite rechnet wie der Browser (16 + Text + 16 px). Probe der
% Textbreiten mit den gemessenen Werten: siehe \typeout POSITIONEN-DESKTOP im Log.
%
% Marken (Schalter \ifmarken, Standard: an; ohne Marken bauen mit
%   xelatex '\def\ohnemarken{}\input{positionen-desktop.tex}'):
% Stil wie info-desktop.tex (weißer Kreis 17 px, Ziffer Inter 600 in #04273C), mittig auf
% der Kopfzeile, 3 px vor dem Spaltennamen: 1 Swap, 2 Total Profit, 3 Margin.
\documentclass[border=0pt]{standalone}
\usepackage{fontspec}
\usepackage{tikz}
\setmainfont{Inter}% Inter 4.0 (Paket fonts-inter), wie bar.tex
\newfontfamily\halbfett{Inter SemiBold}% Ziffern der Marken (wie panel.pdf)

\newlength\px \setlength\px{0.75bp}% 1 CSS-px
\newcommand\schrift{\fontsize{9bp}{10.5bp}\selectfont}% 12 px / line-height 14 px
\definecolor{block}{RGB}{4,39,60}%      .maintable
\definecolor{reiter}{RGB}{0,19,40}%     ul.nav-tabs
\definecolor{aktiv}{RGB}{14,80,114}%    a.nav-link.active; Zeilen mit 10 % Deckung
\definecolor{label}{RGB}{168,187,190}%  th
\definecolor{gewinn}{RGB}{83,188,81}%   color-green-main
\definecolor{verlust}{RGB}{249,64,86}%  color-red-main
\definecolor{markenziffer}{HTML}{04273C}% Ziffer der Marken (wie panel.pdf)

\newif\ifmarken \markentrue
\ifdefined\ohnemarken \markenfalse\fi

\newcommand\breite{956.2}% 10 + 463,2 + 6 + 10 + 6 + 451 + 10
\newcommand\hoehe{145}%   zweite Zeile endet bei 135, + 10 px
\newcommand\links{-141.8}% Verschiebung linker Teil (Instrument beginnt bei x10)
\newcommand\rechts{-597.8}% Verschiebung rechter Teil
\newcommand\kante{479.2}% Mitte der linken Bruchkante (Original-x 621); rechte 10 px weiter

% Textbreiten in px (vor dem Bild messen; in TikZ ist \nullfont aktiv)
\newlength\tw
\makeatletter
\newcommand\breitepx[2]{\settowidth\tw{\schrift #2}%
  \edef#1{\strip@pt\dimexpr\tw*4/3*7200/7227\relax}}
\newcommand\probe[2]{\breitepx\probewert{#1}%
  \typeout{POSITIONEN-DESKTOP: '#1' \probewert\space px (gemessen #2)}}
\makeatother
\probe{Active (0)}{55,0}\probe{Pending (0)}{66,1}\probe{History}{40,6}
\probe{Instrument}{60,7}\probe{Type}{28,1}\probe{Amount}{44,1}
\probe{Trade volume}{78,4}\probe{Trade Volume}{78,4}\probe{Swap}{31,6}
\probe{Commission}{69,8}\probe{Total Profit}{61,5}\probe{Margin}{39,3}
\breitepx\wa{Active}\breitepx\wp{Pending}\breitepx\wh{History}
\makeatother
\pgfmathsetmacro\xh{946.2-(\wh+32)}% Reiter rechtsbündig am Tabellenrand x946,2
\pgfmathsetmacro\xp{\xh-(\wp+32)}
\pgfmathsetmacro\xa{\xp-(\wa+32)}
\typeout{POSITIONEN-DESKTOP: Reiter ab x \xa, \xp, \xh\space bis 946.2 px}

% \flaeche{<Zeilenoberkante>}: getönte Zeile, Zellen-Fläche unter dem 1,6-px-Rand;
% außen Radius 3 px (oben innen 3 x 1,4 px), innen bis über die Bruchkante hinaus
\newcommand\flaeche[1]{%
  \fill[aktiv,fill opacity=0.1] (\breite/2,#1+1.6) -- (13,#1+1.6)
    arc[start angle=270,end angle=180,x radius=3,y radius=1.4] -- (10,#1+26.6)
    arc[start angle=180,end angle=90,radius=3] -- (\breite-13,#1+29.6)
    arc[start angle=90,end angle=0,radius=3] -- (\breite-10,#1+3)
    arc[start angle=0,end angle=-90,x radius=3,y radius=1.4] -- cycle;}
% \hinterteil: Block und Zeilen (werden je Teil einmal geclippt gezeichnet, ohne Text)
\newcommand\hinterteil{%
  \fill[block,rounded corners=7\px] (0,0) rectangle (\breite,\hoehe);
  \flaeche{75.8}\flaeche{105.4}}
% Bruchkante: Zickzack ±3 px, halbe Periode 6 px, über die volle Höhe
\newcommand\zickzack[1]{\foreach \k in {0,...,26} { -- ({#1-3+6*mod(\k,2)},{-6+6*\k}) }}

% \kopf{<Original-x des span>}{<Text>} / \zelle{<Mitte>}{<y>}{<Farbe>}{<Text>}
\newcommand\kopf[2]{\node[anchor=base west,text=label] at (#1,65.4) {#2};}
\newcommand\zelle[4]{\node[anchor=base,text=#3] at (#1,#2) {#4};}

% \marke{<Original-x des Spaltennamens>}{<Nummer>}: Kreis 17 px, Mitte 3 px + Radius davor
\newcommand\marke[2]{%
  \fill[white] (#1-11.5,61) circle[radius=8.5];
  \node[anchor=base,text=markenziffer,font=\fontsize{6.954bp}{8bp}\halbfett]
    at (#1-11.5,61+3.091) {#2};}% 12 px * 17/22

\begin{document}%
\begin{tikzpicture}[x=\px,y={(0,-\px)},font=\schrift,
    every node/.style={inner sep=0,outer sep=0}]
  \path[use as bounding box] (0,0) rectangle (\breite,\hoehe);
  % Fläche: zwei Teile, getrennt durch die Lücke zwischen den Bruchkanten
  \begin{scope}
    \clip (-5,-6) \zickzack{\kante} -- (-5,150) -- cycle;
    \hinterteil
  \end{scope}
  \begin{scope}
    \clip (\breite+5,-6) \zickzack{\kante+10} -- (\breite+5,150) -- cycle;
    \hinterteil
  \end{scope}
  % Reiter (ohne Zähler), bündig mit dem rechten Tabellenrand
  \fill[reiter,rounded corners=3\px] (\xa,10) rectangle (946.2,40);
  \fill[aktiv] (\xp,10) {[rounded corners=3\px] -- (\xa,10) -- (\xa,40)} -- (\xp,40) -- cycle;
  \node[anchor=base,text=white] at ({\xa+(\wa+32)/2},29.2) {Active};
  \node[anchor=base,text=white] at ({\xp+(\wp+32)/2},29.2) {Pending};
  \node[anchor=base,text=white] at ({\xh+(\wh+32)/2},29.2) {History};
  % linker Teil (Original-x)
  \begin{scope}[xshift=\links\px]
    \kopf{177}{Instrument}\kopf{298}{Type}\kopf{383.8}{Amount}\kopf{493.5}{Trade Volume}
    \foreach \y in {95.0,124.6} {
      \zelle{209.4}{\y}{white}{Instrument}% Mitte 214,9 - 5,5 (Pfeil-Bild weggelassen)
      \zelle{319.6}{\y}{white}{Buy}
      \zelle{413.35}{\y}{white}{Menge}
      \zelle{540.25}{\y}{white}{Wert}}
  \end{scope}
  % rechter Teil (Original-x)
  \begin{scope}[xshift=\rechts\px]
    \kopf{1113.6}{Swap}\kopf{1207.4}{Commission}\kopf{1344.2}{Total Profit}\kopf{1467.9}{Margin}
    \foreach \y/\farbe in {95.0/gewinn,124.6/verlust} {
      \zelle{1136.9}{\y}{white}{Betrag}
      \zelle{1249.85}{\y}{white}{—}
      \zelle{1382.5}{\y}{\farbe}{Ergebnis}
      \zelle{1495}{\y}{white}{Betrag}}
    \ifmarken
      \marke{1113.6}{1}
      \marke{1344.2}{2}
      \marke{1467.9}{3}
    \fi
  \end{scope}
\end{tikzpicture}%
\end{document}
'):
% Stil wie info-desktop.tex (weißer Kreis 17 px, Ziffer Inter 600 in #04273C), mittig auf
% der Kopfzeile, 3 px vor dem Spaltennamen: 1 Swap, 2 Total Profit, 3 Margin.
\documentclass[border=0pt]{standalone}
\usepackage{fontspec}
\usepackage{tikz}
\setmainfont{Inter}% Inter 4.0 (Paket fonts-inter), wie bar.tex
\newfontfamily\halbfett{Inter SemiBold}% Ziffern der Marken (wie panel.pdf)

\newlength\px \setlength\px{0.75bp}% 1 CSS-px
\newcommand\schrift{\fontsize{9bp}{10.5bp}\selectfont}% 12 px / line-height 14 px
\definecolor{block}{RGB}{4,39,60}%      .maintable
\definecolor{reiter}{RGB}{0,19,40}%     ul.nav-tabs
\definecolor{aktiv}{RGB}{14,80,114}%    a.nav-link.active; Zeilen mit 10 % Deckung
\definecolor{label}{RGB}{168,187,190}%  th
\definecolor{gewinn}{RGB}{83,188,81}%   color-green-main
\definecolor{verlust}{RGB}{249,64,86}%  color-red-main
\definecolor{markenziffer}{HTML}{04273C}% Ziffer der Marken (wie panel.pdf)

\newif\ifmarken \markentrue
\ifdefined\ohnemarken \markenfalse\fi

\newcommand\breite{956.2}% 10 + 463,2 + 6 + 10 + 6 + 451 + 10
\newcommand\hoehe{145}%   zweite Zeile endet bei 135, + 10 px
\newcommand\links{-141.8}% Verschiebung linker Teil (Instrument beginnt bei x10)
\newcommand\rechts{-597.8}% Verschiebung rechter Teil
\newcommand\kante{479.2}% Mitte der linken Bruchkante (Original-x 621); rechte 10 px weiter

% Textbreiten in px (vor dem Bild messen; in TikZ ist \nullfont aktiv)
\newlength\tw
\makeatletter
\newcommand\breitepx[2]{\settowidth\tw{\schrift #2}%
  \edef#1{\strip@pt\dimexpr\tw*4/3*7200/7227\relax}}
\newcommand\probe[2]{\breitepx\probewert{#1}%
  \typeout{POSITIONEN-DESKTOP: '#1' \probewert\space px (gemessen #2)}}
\makeatother
\probe{Active (0)}{55,0}\probe{Pending (0)}{66,1}\probe{History}{40,6}
\probe{Instrument}{60,7}\probe{Type}{28,1}\probe{Amount}{44,1}
\probe{Trade volume}{78,4}\probe{Trade Volume}{78,4}\probe{Swap}{31,6}
\probe{Commission}{69,8}\probe{Total Profit}{61,5}\probe{Margin}{39,3}
\breitepx\wa{Active}\breitepx\wp{Pending}\breitepx\wh{History}
\makeatother
\pgfmathsetmacro\xh{946.2-(\wh+32)}% Reiter rechtsbündig am Tabellenrand x946,2
\pgfmathsetmacro\xp{\xh-(\wp+32)}
\pgfmathsetmacro\xa{\xp-(\wa+32)}
\typeout{POSITIONEN-DESKTOP: Reiter ab x \xa, \xp, \xh\space bis 946.2 px}

% \flaeche{<Zeilenoberkante>}: getönte Zeile, Zellen-Fläche unter dem 1,6-px-Rand;
% außen Radius 3 px (oben innen 3 x 1,4 px), innen bis über die Bruchkante hinaus
\newcommand\flaeche[1]{%
  \fill[aktiv,fill opacity=0.1] (\breite/2,#1+1.6) -- (13,#1+1.6)
    arc[start angle=270,end angle=180,x radius=3,y radius=1.4] -- (10,#1+26.6)
    arc[start angle=180,end angle=90,radius=3] -- (\breite-13,#1+29.6)
    arc[start angle=90,end angle=0,radius=3] -- (\breite-10,#1+3)
    arc[start angle=0,end angle=-90,x radius=3,y radius=1.4] -- cycle;}
% \hinterteil: Block und Zeilen (werden je Teil einmal geclippt gezeichnet, ohne Text)
\newcommand\hinterteil{%
  \fill[block,rounded corners=7\px] (0,0) rectangle (\breite,\hoehe);
  \flaeche{75.8}\flaeche{105.4}}
% Bruchkante: Zickzack ±3 px, halbe Periode 6 px, über die volle Höhe
\newcommand\zickzack[1]{\foreach \k in {0,...,26} { -- ({#1-3+6*mod(\k,2)},{-6+6*\k}) }}

% \kopf{<Original-x des span>}{<Text>} / \zelle{<Mitte>}{<y>}{<Farbe>}{<Text>}
\newcommand\kopf[2]{\node[anchor=base west,text=label] at (#1,65.4) {#2};}
\newcommand\zelle[4]{\node[anchor=base,text=#3] at (#1,#2) {#4};}

% \marke{<Original-x des Spaltennamens>}{<Nummer>}: Kreis 17 px, Mitte 3 px + Radius davor
\newcommand\marke[2]{%
  \fill[white] (#1-11.5,61) circle[radius=8.5];
  \node[anchor=base,text=markenziffer,font=\fontsize{6.954bp}{8bp}\halbfett]
    at (#1-11.5,61+3.091) {#2};}% 12 px * 17/22

\begin{document}%
\begin{tikzpicture}[x=\px,y={(0,-\px)},font=\schrift,
    every node/.style={inner sep=0,outer sep=0}]
  \path[use as bounding box] (0,0) rectangle (\breite,\hoehe);
  % Fläche: zwei Teile, getrennt durch die Lücke zwischen den Bruchkanten
  \begin{scope}
    \clip (-5,-6) \zickzack{\kante} -- (-5,150) -- cycle;
    \hinterteil
  \end{scope}
  \begin{scope}
    \clip (\breite+5,-6) \zickzack{\kante+10} -- (\breite+5,150) -- cycle;
    \hinterteil
  \end{scope}
  % Reiter (ohne Zähler), bündig mit dem rechten Tabellenrand
  \fill[reiter,rounded corners=3\px] (\xa,10) rectangle (946.2,40);
  \fill[aktiv] (\xp,10) {[rounded corners=3\px] -- (\xa,10) -- (\xa,40)} -- (\xp,40) -- cycle;
  \node[anchor=base,text=white] at ({\xa+(\wa+32)/2},29.2) {Active};
  \node[anchor=base,text=white] at ({\xp+(\wp+32)/2},29.2) {Pending};
  \node[anchor=base,text=white] at ({\xh+(\wh+32)/2},29.2) {History};
  % linker Teil (Original-x)
  \begin{scope}[xshift=\links\px]
    \kopf{177}{Instrument}\kopf{298}{Type}\kopf{383.8}{Amount}\kopf{493.5}{Trade Volume}
    \foreach \y in {95.0,124.6} {
      \zelle{209.4}{\y}{white}{Instrument}% Mitte 214,9 - 5,5 (Pfeil-Bild weggelassen)
      \zelle{319.6}{\y}{white}{Buy}
      \zelle{413.35}{\y}{white}{Menge}
      \zelle{540.25}{\y}{white}{Wert}}
  \end{scope}
  % rechter Teil (Original-x)
  \begin{scope}[xshift=\rechts\px]
    \kopf{1113.6}{Swap}\kopf{1207.4}{Commission}\kopf{1344.2}{Total Profit}\kopf{1467.9}{Margin}
    \foreach \y/\farbe in {95.0/gewinn,124.6/verlust} {
      \zelle{1136.9}{\y}{white}{Betrag}
      \zelle{1249.85}{\y}{white}{—}
      \zelle{1382.5}{\y}{\farbe}{Ergebnis}
      \zelle{1495}{\y}{white}{Betrag}}
    \ifmarken
      \marke{1113.6}{1}
      \marke{1344.2}{2}
      \marke{1467.9}{3}
    \fi
  \end{scope}
\end{tikzpicture}%
\end{document}
'):
% Stil wie info-desktop.tex (weißer Kreis 17 px, Ziffer Inter 600 in #04273C), mittig auf
% der Kopfzeile, 3 px vor dem Spaltennamen: 1 Swap, 2 Total Profit, 3 Margin.
\documentclass[border=0pt]{standalone}
\usepackage{fontspec}
\usepackage{tikz}
\setmainfont{Inter}% Inter 4.0 (Paket fonts-inter), wie bar.tex
\newfontfamily\halbfett{Inter SemiBold}% Ziffern der Marken (wie panel.pdf)

\newlength\px \setlength\px{0.75bp}% 1 CSS-px
\newcommand\schrift{\fontsize{9bp}{10.5bp}\selectfont}% 12 px / line-height 14 px
\definecolor{block}{RGB}{4,39,60}%      .maintable
\definecolor{reiter}{RGB}{0,19,40}%     ul.nav-tabs
\definecolor{aktiv}{RGB}{14,80,114}%    a.nav-link.active; Zeilen mit 10 % Deckung
\definecolor{label}{RGB}{168,187,190}%  th
\definecolor{gewinn}{RGB}{83,188,81}%   color-green-main
\definecolor{verlust}{RGB}{249,64,86}%  color-red-main
\definecolor{markenziffer}{HTML}{04273C}% Ziffer der Marken (wie panel.pdf)

\newif\ifmarken \markentrue
\ifdefined\ohnemarken \markenfalse\fi

\newcommand\breite{956.2}% 10 + 463,2 + 6 + 10 + 6 + 451 + 10
\newcommand\hoehe{145}%   zweite Zeile endet bei 135, + 10 px
\newcommand\links{-141.8}% Verschiebung linker Teil (Instrument beginnt bei x10)
\newcommand\rechts{-597.8}% Verschiebung rechter Teil
\newcommand\kante{479.2}% Mitte der linken Bruchkante (Original-x 621); rechte 10 px weiter

% Textbreiten in px (vor dem Bild messen; in TikZ ist \nullfont aktiv)
\newlength\tw
\makeatletter
\newcommand\breitepx[2]{\settowidth\tw{\schrift #2}%
  \edef#1{\strip@pt\dimexpr\tw*4/3*7200/7227\relax}}
\newcommand\probe[2]{\breitepx\probewert{#1}%
  \typeout{POSITIONEN-DESKTOP: '#1' \probewert\space px (gemessen #2)}}
\makeatother
\probe{Active (0)}{55,0}\probe{Pending (0)}{66,1}\probe{History}{40,6}
\probe{Instrument}{60,7}\probe{Type}{28,1}\probe{Amount}{44,1}
\probe{Trade volume}{78,4}\probe{Trade Volume}{78,4}\probe{Swap}{31,6}
\probe{Commission}{69,8}\probe{Total Profit}{61,5}\probe{Margin}{39,3}
\breitepx\wa{Active}\breitepx\wp{Pending}\breitepx\wh{History}
\makeatother
\pgfmathsetmacro\xh{946.2-(\wh+32)}% Reiter rechtsbündig am Tabellenrand x946,2
\pgfmathsetmacro\xp{\xh-(\wp+32)}
\pgfmathsetmacro\xa{\xp-(\wa+32)}
\typeout{POSITIONEN-DESKTOP: Reiter ab x \xa, \xp, \xh\space bis 946.2 px}

% \flaeche{<Zeilenoberkante>}: getönte Zeile, Zellen-Fläche unter dem 1,6-px-Rand;
% außen Radius 3 px (oben innen 3 x 1,4 px), innen bis über die Bruchkante hinaus
\newcommand\flaeche[1]{%
  \fill[aktiv,fill opacity=0.1] (\breite/2,#1+1.6) -- (13,#1+1.6)
    arc[start angle=270,end angle=180,x radius=3,y radius=1.4] -- (10,#1+26.6)
    arc[start angle=180,end angle=90,radius=3] -- (\breite-13,#1+29.6)
    arc[start angle=90,end angle=0,radius=3] -- (\breite-10,#1+3)
    arc[start angle=0,end angle=-90,x radius=3,y radius=1.4] -- cycle;}
% \hinterteil: Block und Zeilen (werden je Teil einmal geclippt gezeichnet, ohne Text)
\newcommand\hinterteil{%
  \fill[block,rounded corners=7\px] (0,0) rectangle (\breite,\hoehe);
  \flaeche{75.8}\flaeche{105.4}}
% Bruchkante: Zickzack ±3 px, halbe Periode 6 px, über die volle Höhe
\newcommand\zickzack[1]{\foreach \k in {0,...,26} { -- ({#1-3+6*mod(\k,2)},{-6+6*\k}) }}

% \kopf{<Original-x des span>}{<Text>} / \zelle{<Mitte>}{<y>}{<Farbe>}{<Text>}
\newcommand\kopf[2]{\node[anchor=base west,text=label] at (#1,65.4) {#2};}
\newcommand\zelle[4]{\node[anchor=base,text=#3] at (#1,#2) {#4};}

% \marke{<Original-x des Spaltennamens>}{<Nummer>}: Kreis 17 px, Mitte 3 px + Radius davor
\newcommand\marke[2]{%
  \fill[white] (#1-11.5,61) circle[radius=8.5];
  \node[anchor=base,text=markenziffer,font=\fontsize{6.954bp}{8bp}\halbfett]
    at (#1-11.5,61+3.091) {#2};}% 12 px * 17/22

\begin{document}%
\begin{tikzpicture}[x=\px,y={(0,-\px)},font=\schrift,
    every node/.style={inner sep=0,outer sep=0}]
  \path[use as bounding box] (0,0) rectangle (\breite,\hoehe);
  % Fläche: zwei Teile, getrennt durch die Lücke zwischen den Bruchkanten
  \begin{scope}
    \clip (-5,-6) \zickzack{\kante} -- (-5,150) -- cycle;
    \hinterteil
  \end{scope}
  \begin{scope}
    \clip (\breite+5,-6) \zickzack{\kante+10} -- (\breite+5,150) -- cycle;
    \hinterteil
  \end{scope}
  % Reiter (ohne Zähler), bündig mit dem rechten Tabellenrand
  \fill[reiter,rounded corners=3\px] (\xa,10) rectangle (946.2,40);
  \fill[aktiv] (\xp,10) {[rounded corners=3\px] -- (\xa,10) -- (\xa,40)} -- (\xp,40) -- cycle;
  \node[anchor=base,text=white] at ({\xa+(\wa+32)/2},29.2) {Active};
  \node[anchor=base,text=white] at ({\xp+(\wp+32)/2},29.2) {Pending};
  \node[anchor=base,text=white] at ({\xh+(\wh+32)/2},29.2) {History};
  % linker Teil (Original-x)
  \begin{scope}[xshift=\links\px]
    \kopf{177}{Instrument}\kopf{298}{Type}\kopf{383.8}{Amount}\kopf{493.5}{Trade Volume}
    \foreach \y in {95.0,124.6} {
      \zelle{209.4}{\y}{white}{Instrument}% Mitte 214,9 - 5,5 (Pfeil-Bild weggelassen)
      \zelle{319.6}{\y}{white}{Buy}
      \zelle{413.35}{\y}{white}{Menge}
      \zelle{540.25}{\y}{white}{Wert}}
  \end{scope}
  % rechter Teil (Original-x)
  \begin{scope}[xshift=\rechts\px]
    \kopf{1113.6}{Swap}\kopf{1207.4}{Commission}\kopf{1344.2}{Total Profit}\kopf{1467.9}{Margin}
    \foreach \y/\farbe in {95.0/gewinn,124.6/verlust} {
      \zelle{1136.9}{\y}{white}{Betrag}
      \zelle{1249.85}{\y}{white}{—}
      \zelle{1382.5}{\y}{\farbe}{Ergebnis}
      \zelle{1495}{\y}{white}{Betrag}}
    \ifmarken
      \marke{1113.6}{1}
      \marke{1344.2}{2}
      \marke{1467.9}{3}
    \fi
  \end{scope}
\end{tikzpicture}%
\end{document}
'):
% Stil wie info-desktop.tex (weißer Kreis 17 px, Ziffer Inter 600 in #04273C), mittig auf
% der Kopfzeile, 3 px vor dem Spaltennamen: 1 Swap, 2 Total Profit, 3 Margin.
\documentclass[border=0pt]{standalone}
\usepackage{fontspec}
\usepackage{tikz}
\setmainfont{Inter}% Inter 4.0 (Paket fonts-inter), wie bar.tex
\newfontfamily\halbfett{Inter SemiBold}% Ziffern der Marken (wie panel.pdf)

\newlength\px \setlength\px{0.75bp}% 1 CSS-px
\newcommand\schrift{\fontsize{9bp}{10.5bp}\selectfont}% 12 px / line-height 14 px
\definecolor{block}{RGB}{4,39,60}%      .maintable
\definecolor{reiter}{RGB}{0,19,40}%     ul.nav-tabs
\definecolor{aktiv}{RGB}{14,80,114}%    a.nav-link.active; Zeilen mit 10 % Deckung
\definecolor{label}{RGB}{168,187,190}%  th
\definecolor{gewinn}{RGB}{83,188,81}%   color-green-main
\definecolor{verlust}{RGB}{249,64,86}%  color-red-main
\definecolor{markenziffer}{HTML}{04273C}% Ziffer der Marken (wie panel.pdf)

\newif\ifmarken \markentrue
\ifdefined\ohnemarken \markenfalse\fi

\newcommand\breite{956.2}% 10 + 463,2 + 6 + 10 + 6 + 451 + 10
\newcommand\hoehe{145}%   zweite Zeile endet bei 135, + 10 px
\newcommand\links{-141.8}% Verschiebung linker Teil (Instrument beginnt bei x10)
\newcommand\rechts{-597.8}% Verschiebung rechter Teil
\newcommand\kante{479.2}% Mitte der linken Bruchkante (Original-x 621); rechte 10 px weiter

% Textbreiten in px (vor dem Bild messen; in TikZ ist \nullfont aktiv)
\newlength\tw
\makeatletter
\newcommand\breitepx[2]{\settowidth\tw{\schrift #2}%
  \edef#1{\strip@pt\dimexpr\tw*4/3*7200/7227\relax}}
\newcommand\probe[2]{\breitepx\probewert{#1}%
  \typeout{POSITIONEN-DESKTOP: '#1' \probewert\space px (gemessen #2)}}
\makeatother
\probe{Active (0)}{55,0}\probe{Pending (0)}{66,1}\probe{History}{40,6}
\probe{Instrument}{60,7}\probe{Type}{28,1}\probe{Amount}{44,1}
\probe{Trade volume}{78,4}\probe{Trade Volume}{78,4}\probe{Swap}{31,6}
\probe{Commission}{69,8}\probe{Total Profit}{61,5}\probe{Margin}{39,3}
\breitepx\wa{Active}\breitepx\wp{Pending}\breitepx\wh{History}
\makeatother
\pgfmathsetmacro\xh{946.2-(\wh+32)}% Reiter rechtsbündig am Tabellenrand x946,2
\pgfmathsetmacro\xp{\xh-(\wp+32)}
\pgfmathsetmacro\xa{\xp-(\wa+32)}
\typeout{POSITIONEN-DESKTOP: Reiter ab x \xa, \xp, \xh\space bis 946.2 px}

% \flaeche{<Zeilenoberkante>}: getönte Zeile, Zellen-Fläche unter dem 1,6-px-Rand;
% außen Radius 3 px (oben innen 3 x 1,4 px), innen bis über die Bruchkante hinaus
\newcommand\flaeche[1]{%
  \fill[aktiv,fill opacity=0.1] (\breite/2,#1+1.6) -- (13,#1+1.6)
    arc[start angle=270,end angle=180,x radius=3,y radius=1.4] -- (10,#1+26.6)
    arc[start angle=180,end angle=90,radius=3] -- (\breite-13,#1+29.6)
    arc[start angle=90,end angle=0,radius=3] -- (\breite-10,#1+3)
    arc[start angle=0,end angle=-90,x radius=3,y radius=1.4] -- cycle;}
% \hinterteil: Block und Zeilen (werden je Teil einmal geclippt gezeichnet, ohne Text)
\newcommand\hinterteil{%
  \fill[block,rounded corners=7\px] (0,0) rectangle (\breite,\hoehe);
  \flaeche{75.8}\flaeche{105.4}}
% Bruchkante: Zickzack ±3 px, halbe Periode 6 px, über die volle Höhe
\newcommand\zickzack[1]{\foreach \k in {0,...,26} { -- ({#1-3+6*mod(\k,2)},{-6+6*\k}) }}

% \kopf{<Original-x des span>}{<Text>} / \zelle{<Mitte>}{<y>}{<Farbe>}{<Text>}
\newcommand\kopf[2]{\node[anchor=base west,text=label] at (#1,65.4) {#2};}
\newcommand\zelle[4]{\node[anchor=base,text=#3] at (#1,#2) {#4};}

% \marke{<Original-x des Spaltennamens>}{<Nummer>}: Kreis 17 px, Mitte 3 px + Radius davor
\newcommand\marke[2]{%
  \fill[white] (#1-11.5,61) circle[radius=8.5];
  \node[anchor=base,text=markenziffer,font=\fontsize{6.954bp}{8bp}\halbfett]
    at (#1-11.5,61+3.091) {#2};}% 12 px * 17/22

\begin{document}%
\begin{tikzpicture}[x=\px,y={(0,-\px)},font=\schrift,
    every node/.style={inner sep=0,outer sep=0}]
  \path[use as bounding box] (0,0) rectangle (\breite,\hoehe);
  % Fläche: zwei Teile, getrennt durch die Lücke zwischen den Bruchkanten
  \begin{scope}
    \clip (-5,-6) \zickzack{\kante} -- (-5,150) -- cycle;
    \hinterteil
  \end{scope}
  \begin{scope}
    \clip (\breite+5,-6) \zickzack{\kante+10} -- (\breite+5,150) -- cycle;
    \hinterteil
  \end{scope}
  % Reiter (ohne Zähler), bündig mit dem rechten Tabellenrand
  \fill[reiter,rounded corners=3\px] (\xa,10) rectangle (946.2,40);
  \fill[aktiv] (\xp,10) {[rounded corners=3\px] -- (\xa,10) -- (\xa,40)} -- (\xp,40) -- cycle;
  \node[anchor=base,text=white] at ({\xa+(\wa+32)/2},29.2) {Active};
  \node[anchor=base,text=white] at ({\xp+(\wp+32)/2},29.2) {Pending};
  \node[anchor=base,text=white] at ({\xh+(\wh+32)/2},29.2) {History};
  % linker Teil (Original-x)
  \begin{scope}[xshift=\links\px]
    \kopf{177}{Instrument}\kopf{298}{Type}\kopf{383.8}{Amount}\kopf{493.5}{Trade Volume}
    \foreach \y in {95.0,124.6} {
      \zelle{209.4}{\y}{white}{Instrument}% Mitte 214,9 - 5,5 (Pfeil-Bild weggelassen)
      \zelle{319.6}{\y}{white}{Buy}
      \zelle{413.35}{\y}{white}{Menge}
      \zelle{540.25}{\y}{white}{Wert}}
  \end{scope}
  % rechter Teil (Original-x)
  \begin{scope}[xshift=\rechts\px]
    \kopf{1113.6}{Swap}\kopf{1207.4}{Commission}\kopf{1344.2}{Total Profit}\kopf{1467.9}{Margin}
    \foreach \y/\farbe in {95.0/gewinn,124.6/verlust} {
      \zelle{1136.9}{\y}{white}{Betrag}
      \zelle{1249.85}{\y}{white}{—}
      \zelle{1382.5}{\y}{\farbe}{Ergebnis}
      \zelle{1495}{\y}{white}{Betrag}}
    \ifmarken
      \marke{1113.6}{1}
      \marke{1344.2}{2}
      \marke{1467.9}{3}
    \fi
  \end{scope}
\end{tikzpicture}%
\end{document}
